\documentclass[10pt,a4paper]{amsart}
\usepackage[margin=19mm]{geometry}
\usepackage{amsmath,amssymb,mathtools}
\usepackage[scr=rsfs,cal=euler]{mathalfa}
\usepackage{xfrac}
\usepackage[unicode,hidelinks,bookmarks=false]{hyperref}
\newcommand{\ie}{i.e.\ }
\newcommand{\cf}{cf.\ }
\newcommand{\resp}{respectively\ }
\newcommand{\Subsection}{\subsection}
\providecommand{\nobreakdash}{\nobreak\hbox{-}}
\setlength{\emergencystretch}{2em}
\multlinegap0pt
\usepackage{amssymb}
\newtheorem{theorem}{Theorem}
\newcommand{\Z}{\mathbb Z}
\newcommand{\vTwo}{\nu_2}
\title{An infinite small-step $\Z^3$-walk with no collinear triple}
\author{Stijn Cambie, Erik Kalviainen}
\date{Source: arXiv:2609.01766v1 (CC BY 4.0)}
\begin{document}
\maketitle

\noindent\emph{Source and licence.} An infinite small-step $\Z^3$-walk with no collinear triple. Version arXiv:2609.01766v1 (CC BY 4.0), distributed by arXiv under the Creative Commons Attribution 4.0 International licence, http://creativecommons.org/licenses/by/4.0/, as recorded in the arXiv licence metadata for this exact version and shown on the abstract page https://arxiv.org/abs/2609.01766v1. Full licence notice: \texttt{licenses/arxiv-2609.01766-CC-BY-4.0.txt}.\par\smallskip

\noindent\textit{Editorial scope.} Counted unit: the article's theorem theorem (source0.tex lines 28--36). The article's complete original proof of it is delivered with it (source0.tex lines 38--128, one \texttt{proof} environment of 91 lines). No mathematics is written, repaired or re-derived: the statement and the proof are the article's own lines, in the article's own notation, and no retained line is substituted or re-flowed.  No further source-local context is needed: every cross-reference of every retained line resolves inside the two delivered spans.  Nothing was omitted from any retained span, so the package carries no omission record.  No retained line contains a citation, so the wrapper carries no bibliography and no bibliography entry.  No retained line contains a figure environment, an image-inclusion command or drawing code, so no image is delivered and the package declares no asset dependency.  Editorial formatting only, in addition: an AMS \texttt{amsart} preamble that restates the article's own declarations and macros used by the retained lines, namely the environments \texttt{theorem} on separate counters, together with the standard packages those lines need.  The retained environments print in this document as Theorem 1 in order of appearance, whereas the article prints the same items with the numbering of its own full document.  The complete native source of the article as distributed in the arXiv e-print for this exact version is bundled unchanged as \texttt{sources/209141/source0.tex}.
\begin{theorem}
There is an infinite sequence $P_0,P_1,\ldots$ in $\Z^3$ with no three
collinear terms and with
\[
 P_{n+1}-P_n\in\{-2,-1,0,1,2\}^2\times\{1,2,\ldots,7\}
 \qquad(n\geq0).
\]
Only sixteen successive displacement vectors occur.
\end{theorem}

\begin{proof}
Let $s_2(n)$ be the number of $1$s in the binary expansion of $n$, and put
\[
 u_n=i^{s_2(n)},\qquad z_n=\sum_{0\leq r<n}u_r\in\Z[i].
\]
Thus $(z_n)$ is a nearest-neighbor walk in the Gaussian lattice.  Binary
expansion gives, for $\varepsilon\in\{0,1\}$,
\begin{equation}\label{eq:recurrence}
 u_{2n+\varepsilon}=i^\varepsilon u_n,
 \qquad z_{2n+\varepsilon}=(1+i)z_n+\varepsilon u_n.
\end{equation}

We first record the only property of this walk that we need.  If $0\leq m<n$
and $u_m=u_n$, then
\begin{equation}\label{eq:same-state}
 \vTwo\bigl(|z_n-z_m|^2\bigr)=\vTwo(n-m).
\end{equation}
Indeed, whenever $n-m$ is even, the endpoints have the same parity, say
$m=2a+\varepsilon$ and $n=2b+\varepsilon$.  From
\eqref{eq:recurrence}, $u_m=u_n$ implies $u_a=u_b$ and
\[
 z_n-z_m=(1+i)(z_b-z_a).
\]
If $r=\vTwo(n-m)$, repeating this reduction $r$ times gives
\[
 z_n-z_m=(1+i)^r(z_{n'}-z_{m'}),
\]
where $n'-m'$ is odd.  The last difference is a sum of an odd number of
Gaussian units.  Writing it as $x+iy$, we have $x+y$ odd, so $x^2+y^2$ is
odd.  Since $|1+i|^2=2$, equation~\eqref{eq:same-state} follows.

Let $\alpha_n\in\{0,1,2,3\}$ be determined by $u_n=i^{\alpha_n}$.  Tag the
four states by the corners
\[
 c_0=0,\qquad c_1=i,\qquad c_2=-1+i,\qquad c_3=-1
\]
of a unit square in cyclic order, and define
\begin{equation}\label{eq:points}
 w_n=2z_n+c_{\alpha_n},\qquad
 h_n=4n+\alpha_n,\qquad
 P_n=(\Re w_n,\Im w_n,h_n).
\end{equation}
We claim that, for every $m<n$,
\begin{equation}\label{eq:all-pairs}
 \vTwo\bigl(|w_n-w_m|^2\bigr)=\vTwo(h_n-h_m).
\end{equation}
Put $a=\alpha_m$, $b=\alpha_n$, and $d=n-m$.  Then
\[
 w_n-w_m=2(z_n-z_m)+c_b-c_a,
 \qquad h_n-h_m=4d+b-a.
\]
If $a=b$, equation~\eqref{eq:all-pairs} follows from
\eqref{eq:same-state}.  If $b-a$ is odd, exactly one coordinate of
$c_b-c_a$ is odd; hence the squared modulus and $4d+b-a$ are both odd.  If
$b-a=\pm2$, both planar coordinates are odd, so their squared sum is $2$
modulo $4$, while $4d\pm2$ has valuation one.  This proves
\eqref{eq:all-pairs}.

The sequence has bounded steps.  If $j=\alpha_n$ and $k=\alpha_{n+1}$, then
\begin{equation}\label{eq:steps}
 w_{n+1}-w_n=2i^j+c_k-c_j,
 \qquad h_{n+1}-h_n=4+k-j.
\end{equation}
The second quantity lies between $1$ and $7$.  The corner $c_j$ was oriented
so that $i^j$ points outward from the square.  The component of $c_k-c_j$ in
that direction is $0$ or $-1$, and its perpendicular component has absolute
value at most $1$.  Thus both planar step coordinates have absolute value at
most $2$.  Equation~\eqref{eq:steps} also shows that the step is determined by
$(j,k)$, so at most sixteen steps occur.  Since $h_{n+1}>h_n$, the points are
distinct.

It remains to exclude collinearity.  Suppose $P_a,P_b,P_c$ are collinear for
$a<b<c$, and write
\[
 A=h_b-h_a,\qquad B=h_c-h_b,
 \qquad X=w_b-w_a,\qquad Y=w_c-w_b.
\]
Because $A,B>0$, collinearity gives one common complex slope,
\[
 \frac XA=\frac YB=\frac{X+Y}{A+B}.
\]
Extend $\vTwo$ to nonzero rational numbers in the usual way.  Applying
\eqref{eq:all-pairs} to the three chords and taking squared moduli of the
slopes yields
\[
 \vTwo(A)=\vTwo(B)=\vTwo(A+B).
\]
This is impossible: if the first two valuations equal $t$, then $A/2^t$ and
$B/2^t$ are odd, so $(A+B)/2^t$ is even.  Hence no three points $P_n$ are
collinear.
\end{proof}

\end{document}
